% 06_tablas/tabla_analisis_documental.tex
% Tabla 2 — Matriz de análisis documental y legislativo (Cap. III, sec. 3.6)

\begingroup
\small
\setlength{\LTpre}{6pt}
\setlength{\LTpost}{6pt}

\begin{longtable}{>{\raggedright}p{0.025\textwidth}%
                  >{\raggedright}p{0.11\textwidth}%
                  >{\raggedright}p{0.31\textwidth}%
                  >{\raggedright}p{0.05\textwidth}%
                  >{\raggedright}p{0.22\textwidth}%
                  >{\raggedright\arraybackslash}p{0.14\textwidth}}

\caption{Matriz de análisis documental y legislativo de la investigación}
\label{tab:analisis_documental}\\

\toprule
\textbf{N.°} &
\textbf{Tipo} &
\textbf{Documento} &
\textbf{Año} &
\textbf{Estado jurídico / documental} &
\textbf{OE principal} \\
\midrule
\endfirsthead

\multicolumn{6}{l}{\small\textit{Tabla~\ref{tab:analisis_documental} (continuación)}}\\
\toprule
\textbf{N.°} &
\textbf{Tipo} &
\textbf{Documento} &
\textbf{Año} &
\textbf{Estado jurídico / documental} &
\textbf{OE principal} \\
\midrule
\endhead

\midrule
\multicolumn{6}{r}{\small\textit{Continúa en la página siguiente}}\\
\endfoot

\bottomrule
\endlastfoot

% ------- NORMATIVA -------
1 & Normativa constitucional & Constitución de la República del Ecuador & 2008 & Vigente. Base constitucional del sistema de protección ambiental y derechos de la naturaleza (arts.\ 14, 71--74, 316, 395--399, 407--409) & OE1, OE4 \\
\addlinespace
2 & Normativa legal & Código Orgánico del Ambiente & 2017 & Vigente. Principal instrumento codificador del derecho ambiental ecuatoriano. Define licenciamiento, prevención y principios ambientales & OE1, OE3 \\
\addlinespace
3 & Normativa reglamentaria & Reglamento al Código Orgánico del Ambiente & 2019 & Vigente (Decreto Ejecutivo 752). Regula licenciamiento, consulta ambiental, categorización SUIA. Arts.\ 462 y 463 declarados inconstitucionales & OE1, OE3 \\
\addlinespace
4 & Normativa legal & Ley de Minería (texto consolidado hasta agosto de 2025) & 2009/2025 & Vigente con reformas sucesivas. Referencia comparativa para identificar cambios introducidos por la Ley de Fortalecimiento 2026 & OE1, OE2 \\
\addlinespace
5 & Normativa legal & Ley Orgánica Reformatoria a la Ley de Minería & 2013 & Vigente (incorporada al texto de la Ley de Minería). Antecedente legislativo de la tendencia reformadora del sector minero & OE2 \\
\addlinespace
6 & Normativa legal & Ley Orgánica para el Fortalecimiento de los Sectores Estratégicos de Minería y Energía & 2026 & Vigente (R.O.\ Quinto Suplemento 234, 2-mar-2026). Al 5-abr-2026 acumula al menos 11 demandas de inconstitucionalidad ante la CCE & OE2, OE3 \\
\addlinespace
7 & Normativa reglamentaria & Reglamento General a la Ley de Minería & 2009 & Vigente. Regula el ciclo minero: exploración, explotación, beneficio, cierre. Modificado parcialmente en 2019 y 2025 & OE1 \\
\addlinespace
8 & Normativa legal & Código Orgánico Integral Penal & 2014 & Vigente. Arts.\ 260 y 261: tipos penales de delitos contra los recursos naturales (minería ilegal, daño ambiental) & OE1 \\
\midrule
% ------- JURISPRUDENCIA -------
9 & Jurisprudencia constitucional & Sentencia No.\ 1149-19-JP/21 — Caso Bosque Protector Los Cedros & 2021 & Precedente vinculante. Establece estándares de consulta ambiental (art.\ 398 CRE), principio de precaución e \textit{in dubio pro natura} en concesiones mineras & OE1, OE3 \\
\addlinespace
10 & Jurisprudencia constitucional & Sentencia No.\ 32-17-IN/21 — Inconstitucionalidad arts.\ 86 y 136 RAAM & 2021 & Precedente vinculante. Establece reserva de ley orgánica como límite a normas reglamentarias con incidencia en derechos constitucionales & OE1, OE2 \\
\addlinespace
11 & Jurisprudencia constitucional & Sentencia No.\ 22-18-IN/21 — Derechos de la naturaleza; arts.\ 462-463 RCOAM & 2021 & Precedente vinculante. Distingue consulta ambiental y CPLI. Declara inconstitucionales arts.\ 462 y 463 RCOAM & OE1, OE3 \\
\addlinespace
12 & Documento institucional jurisprudencial & Guía de jurisprudencia constitucional sobre derechos de la naturaleza (CEDEC) & 2023 & Documento de sistematización. No es sentencia autónoma. Compila 10 sentencias sobre DDN con valor de referencia sistémica & OE1 \\
\midrule
% ------- INTERNACIONAL -------
13 & Instrumento internacional & Acuerdo Regional de Escazú & 2018 & Tratado en vigor para Ecuador (22-abr-2021). Establece derechos de acceso: información, participación y justicia en asuntos ambientales & OE1, OE3 \\
\addlinespace
14 & Instrumento internacional & CEPAL / MAATE — Ruta para la implementación del Acuerdo de Escazú en Ecuador & 2023 & Documento técnico no vinculante. Identifica brechas normativas del ordenamiento ecuatoriano frente a los estándares de Escazú & OE3 \\
\midrule
% ------- INSTITUCIONAL -------
15 & Documento institucional & ARCOM — Informe de Rendición de Cuentas 2024 & 2024 & Documento oficial. Datos de fiscalización (57\% de meta), 387 puntos de minería ilegal activos, estructura institucional & OE3 \\
\midrule
% ------- DOCTRINA -------
16 & Doctrina académica & Grefa Valencia — Derechos de la naturaleza y extractivismo minero en Ecuador (tesis maestría UASB) & 2021 & Tesis de maestría verificada. Análisis de la tensión DDN-extractivismo con enfoque constitucional & OE1, OE3 \\
\addlinespace
17 & Doctrina académica & Villacís Calvas — Seguridad jurídica e irretroactividad en concesiones de pequeña minería (tesis maestría UASB) & 2022 & Tesis de maestría verificada. Desarrollo dogmático del concepto de seguridad jurídica aplicado al sector minero & OE1, OE2 \\
\addlinespace
18 & Doctrina académica & Lozano Espinosa — La naturaleza como sujeto de derechos desde la perspectiva de la CCE (2018--2022) & 2023 & Artículo científico indexado (\textit{Multiverso Journal}). Sistematiza la jurisprudencia constitucional sobre DDN & OE1 \\
\addlinespace
19 & Doctrina académica & Ochoa Jiménez — Los derechos de la naturaleza en Ecuador: análisis desde la jurisprudencia constitucional & 2024 & Artículo científico indexado (\textit{Sur Academia}). Evaluación crítica de la evolución jurisprudencial & OE1 \\
\addlinespace
20 & Doctrina académica & Peck et al.\ — Conflicto entre derechos de la naturaleza y minería en Ecuador: caso Los Cedros & 2024 & Artículo científico indexado (\textit{People and Nature}). Implicaciones del caso Los Cedros para la conservación de biodiversidad & OE1, OE3 \\
\addlinespace
21 & Doctrina académica & Sánchez-Romero, Jumbo-Cañar y Ramírez-López — Minería en áreas protegidas del Ecuador & 2026 & Artículo científico indexado (\textit{RTETS}). Análisis de límites jurídicos e impacto socioambiental & OE3 \\
\addlinespace
22 & Doctrina académica & Medina Llerena y García Erazo — Responsabilidad estatal por omisión en el control de la minería ilegal en la Amazonía & 2026 & Artículo científico indexado (\textit{Polo del Conocimiento}). Tipología de la omisión estatal en fiscalización ambiental & OE3 \\
\midrule
% ------- PRENSA -------
23 & Prensa contextual & 16 fuentes de prensa ecuatoriana (Primicias, El Comercio, GK, Plan V, entre otros) & 2025--2026 & Fuentes de contexto. No reemplazan normas ni sentencias. Se usan para describir debate público, conflictividad y estado de demandas ante la CCE & OE2 (contextual) \\
\midrule
% ------- REFORMAS LEGISLATIVAS -------
24 & Documento legislativo contextual & Noticia institucional AN: aprobación del informe de segundo debate del proyecto de ley del sector minero y energético (24-feb-2026) & 2026 & Fuente contextual. No es texto oficial del informe de segundo debate. No reemplaza el expediente legislativo & OE2 (contextual) \\

\end{longtable}

\endgroup

\vspace{4pt}
\noindent\textit{Nota.} La tabla presenta el corpus documental analizado en la investigación, organizado por tipo de fuente. La columna ``OE principal'' indica el objetivo específico al que contribuye de manera primaria cada fuente, sin excluir aportes secundarios a otros objetivos. Las fuentes de prensa (fila 23) y el documento legislativo contextual (fila 24) no tienen valor jurídico autónomo y se utilizan únicamente para contextualizar el debate público y el estado del trámite legislativo. El estado jurídico de la Ley de Fortalecimiento 2026 (fila 6) reconoce su vigencia formal sin prejuzgar el resultado de los procesos constitucionales en curso. Elaboración propia a partir del corpus documental del repositorio de la investigación.
